% =============================================================
% Section 7. Numerical Verification
%
% Source: A-route-ch3-three-instances-v0.2.md §5 + code/A_route_ch3_numerical.py
% Template: paper/v0.1/notes/narrative.md §4 §7
% =============================================================

\noindent
\cref{thm:strict-equality} reduces $\sheafBW(\samps)$ to a spectral
ratio on the sampling operator $A$. We verify this strict equality
numerically across the three instances of \cref{sec:three_instances}.
The verification confirms the linear-algebraic identity and the code
implementation; it does not constitute independent physical evidence
for each instance, as we make explicit in
\cref{subsec:7-honest-statement}.

% -------------------------------------------------------------
\subsection{Experimental setup}\label{subsec:7-setup}

A self-contained Python script (\texttt{A\_route\_ch3\_numerical.py},
$147$~lines, depending only on
\texttt{numpy} \cite{HarrisEtAl2020NumPy}) constructs a toy operator
$A \in \C^{M \times K}$ for each of the three instances and computes
two quantities for comparison:
\begin{enumerate}[leftmargin=*,label=(\roman*)]
\item the spectral-gap form
  $\sqrt{\gamma_+(\delta^0)/\Gamma_+(\delta^0)}$, where
  $(\delta^0)^* \delta^0 = A^* A$ on the relevant $V_?$ subspace;
\item the singular-value form
  $\sigma_{\min}(A)/\sigma_{\max}(A)$.
\end{enumerate}
For each instance, $A$ is generated by drawing a random orthonormal
basis $U \in \C^{N \times K}$ via QR factorization of a complex
Gaussian matrix and then sub-sampling $M$ rows of $U$ to obtain
$A$, mirroring the row-sampled basis representation
$A_{n,l} = v_l(i_n)$ used uniformly across the three instances of
\cref{sec:three_instances}. The dimensions are
\begin{center}
\begin{tabular}{lcccl}
\toprule
Instance & $N$ & $K$ & $M$ & Subspace label \\
\midrule
1: TopSP & $100$ & $10$ & $30$ & $V_B^{\TopSP}$ \\
2: FEM--EM & $50$ & $6$ & $8$ & $V_S^{\FEM}$ \\
3: BZ & $30$ & $4$ & $12$ & $V_{\occ}^{\BZ}$ \\
\bottomrule
\end{tabular}
\end{center}
The random seed is fixed at $\mathtt{seed} = 20260507$ for
reproducibility. The relative deviation between the two forms is
defined as
\[
  \mathrm{reldiff}
  \;:=\;
  \frac{\bigl|\sqrt{\gamma_+/\Gamma_+} - \sigma_{\min}/\sigma_{\max}\bigr|}{\sigma_{\min}/\sigma_{\max}}.
\]

% -------------------------------------------------------------
\subsection{Results}\label{subsec:7-results}

\begin{table}[h]
\centering
\begin{tabular}{lccc}
\toprule
\textbf{Instance} & $\sqrt{\gamma_+/\Gamma_+}$ & $\sigma_{\min}/\sigma_{\max}$ & $\mathrm{reldiff}$ \\
\midrule
1: TopSP   & $0.367756144311$ & $0.367756144311$ & $7.55 \times 10^{-16}$ \\
2: FEM--EM & $0.053301378959$ & $0.053301378959$ & $1.64 \times 10^{-14}$ \\
3: BZ      & $0.519194438983$ & $0.519194438983$ & $2.14 \times 10^{-16}$ \\
\bottomrule
\end{tabular}
\caption{Numerical verification of \cref{thm:strict-equality} on the
  three instances at $\mathtt{seed} = 20260507$. The two forms
  (spectral-gap and singular-value) of \cref{def:sheafBW} agree to at
  least twelve decimal digits in every instance. Instances~1 and~3
  attain $\mathrm{reldiff}$ within $4 \times \epsilon_{\mathrm{mach}}$
  (IEEE~754 unit roundoff $\epsilon_{\mathrm{mach}} \approx
  2.22 \times 10^{-16}$); Instance~2 deviates by $1.64 \times 10^{-14}$,
  reflecting the expected roundoff amplification of comparing
  $\sqrt{\lambda_{\min}(A^* A)/\lambda_{\max}(A^* A)}$ with
  $\sigma_{\min}(A)/\sigma_{\max}(A)$ when $A$ is poorly conditioned,
  since $\kappa(A^* A) = \kappa(A)^2$. All three deviations are at
  least an order of magnitude below the floating-point error budget for
  this comparison on matrices of the given shapes.}
\label{tab:7-results}
\end{table}

The maximum relative deviation across the three instances is
$1.64 \times 10^{-14}$ (Instance~2), driven by the squaring of the
condition number under the $A^* A$ form rather than by any numerical
failure of \cref{thm:strict-equality}. The spectral-gap form and the
singular-value form agree to within the floating-point error budget
in every instance.

% -------------------------------------------------------------
\subsection{Honest statement on the nature of this verification}\label{subsec:7-honest-statement}

The script of \cref{subsec:7-setup} verifies the spectral-gap form
$\sqrt{\gamma_+/\Gamma_+}$ and the singular-value form
$\sigma_{\min}/\sigma_{\max}$ of \cref{def:sheafBW} agree numerically
on a randomly generated $A$ of the relevant shape. Two facts are
verified:
\begin{itemize}[leftmargin=*]
\item the linear-algebraic identity
  $\sqrt{\lambda_{\min}(A^* A) / \lambda_{\max}(A^* A)}
  = \sigma_{\min}(A)/\sigma_{\max}(A)$,
  which holds for any $M \times K$ matrix $A$ (a standard consequence
  of the singular value decomposition);
\item the code implementation (\texttt{A\_route\_ch3\_numerical.py})
  computes both forms without introducing numerical artifacts beyond
  IEEE 754 roundoff.
\end{itemize}
What the script does \emph{not} do:
\begin{itemize}[leftmargin=*]
\item It does not provide independent physical evidence that the
  domain $V_?$ in each instance is correctly chosen for the underlying
  application (signal reconstruction in TopSP, finite element solution
  in FEM--EM, band-structure integration in BZ). All three instances
  use the same toy construction
  $U = \mathrm{QR}(\mathrm{randn}(N, K))$, which only matches the
  abstract structural template
  ``finite-dim Hilbert subspace + row-sampled basis representation.''
\item It does not validate the heuristic skeleton of
  \cref{empirical:8b}: that proposition rests on the crossed
  evidence channels of \cref{tab:6-evidence}, not on this script.
\item It does not exercise the cellular-sheaf construction of
  \cref{sec:sheafbw} independently: by \cref{prop:delta-equals-A},
  $\delta^0 = A$ on the sampling sheaf, so
  $(\delta^0)^* \delta^0 = A^* A$ holds by definition rather than as
  a numerical fact. The deviation reported in \cref{tab:7-results}
  is the IEEE 754 roundoff of computing eigenvalues of $A^*A$ versus
  squared singular values of $A$; it does not test the cellular
  sheaf structure.
\end{itemize}
The script is best understood as a sanity check on the algebra of
\cref{def:sheafBW} and on the code path that computes it, not as
empirical support for the framework's applicability to any specific
domain.

% -------------------------------------------------------------
\subsection{Physical evidence per instance}\label{subsec:7-evidence}

The empirical content of each instance rests on the following sources,
which are not exercised by the script of \cref{subsec:7-setup}.

\begin{table}[h]
\centering
\small
\begin{tabular}{@{}l p{0.45\linewidth} p{0.32\linewidth}@{}}
\toprule
\textbf{Instance} & \textbf{Empirical content} & \textbf{Source} \\
\midrule
1: TopSP    & $35$ MRI sampling configurations; full column rank in all $35$;
              Cartesian $\sheafBW \approx 0.0091$ vs.\ radial $\approx 2.07 \times 10^{-5}$
            & this paper, project simulation suite (\cref{subsec:6-topsp}) \\
2: FEM--EM  & finite-dimensional restriction is a project-internal scope
              statement, not an empirical claim
            & \cref{subsec:6-femem} \\
3: BZ       & five decades of condensed-matter practice consistently report
              non-Cartesian tessellations outperforming Cartesian quadrature
            & \cite{Blochl1994,LehmannTaut1972,JepsenAndersen1971,MonkhorstPack1976} \\
\bottomrule
\end{tabular}
\caption{Sources of empirical content per instance. Instance 2 is
  established by the finite-dimensional restriction of
  \cref{subsec:sheafbw-hilbert-complex}; the row labeled ``finite-dimensional
  restriction'' is a project-internal scope statement and is not a
  software- or hardware-derived empirical claim.}
\label{tab:7-evidence}
\end{table}

\noindent
The numerical verification closes the main argument of the paper. We
conclude with a discussion of the two independent open problems left
by this work.
